\sscp{\tsc{\ili{Kisar}-Luang}}{04-06-02}
A node containing \ili{Roma},
\ili{Kisar}, and the \ili{Luang} cluster is well defined on the basis of
*z > **ty [ʧ] which is not only unique within the region,
but almost unique among all Malayo-Polynesian languages.\footnote{
		The nearest languages with *z > \it{ʧ} are in Borneo \citep[279]{Smith2017a}.}
We call this node \tsc{\ili{Kisar}-Luang} after the two most socially
prominent islands on which these languages are spoken.
\ili{As} discussed above, the change *z > **ty [ʧ] may provide evidence
for a more extended node within Southwest Maluku also including \tsc{\ili{Teun}-\ili{Nila}-Serua} and \tsc{Babar}.
Additional evidence for \tsc{\ili{Kisar}-Luang} comes from shared *s
> \it{h, s}; *j > \it{r}; *aw > \it{a};
and *ŋ > \it{n} in most positions,
of which *j > \it{r}
and *-aw > \it{a} also occur widely among \tsc{Southwest Maluku} languages.

\begin{table}\tcp{\tsc{\ili{Kisar}-Luang} *z}{04-27}
\begin{adjustbox}{width=\textwidth}
	\begin{threeparttable}\st{2pt}
	\begin{tabular}{lll@{ }lllll}\lsptoprule
*gloss	&	path	&	far	&	Job’s tears	&	rain	&	point	&	ladder	&	green\\
PMP	&	*\tbr{z}alan	&	*\tbr{z}auq	&	*\tbr{z}əlay\su{†}	&	*qu\tbr{z}an	&	*tu\tbr{z}uq	&	*haRə\tbr{z}an	&	**mo\tbr{z}o\su{‡}\\ \midrule
\ili{Roma}	&	{\hr}\tbr{ty}alla	&	{\hr}a\tbr{ty}ounna\su{Δ}	&	{\hr}\tbr{t}eli	&	{\hr}u\tbr{t}ni	&		&	{\hr}re\tbr{t}ni	&	{\hr}mo{\tl}mo\tbr{t}i\\
\ili{Kisar}	&	{\hr}\tbr{k}alla	&	{\hr}\tbr{k}oʔu	&	{\hr}\tbr{k}eleʔuk	&	{\hr}o\tbr{k}on	&	{\hr}ku\tbr{k}ul	&	{\hr}ro\tbr{k}on	&	{\hr}mok{\tl}mo\tbr{k}e\\
\ili{Wetan} (T)	&	{\hr}\tbr{t}alna	&	\cg{(olieta)}	&		&	{\hr}o\tbr{t}na	&	{\hr}tu\tbr{t}ti	&		&	\\
\ili{Wetan} (W)	&	{\hr}\tbr{t}alla	&	\cg{(oilieta)}	&		&	{\hr}o\tbr{t}na	&	{\hr}tu\tbr{t}i	&	{\hr}re\tbr{t}na	&	{\hr}mo\tbr{t}a\\
\ili{Luang}	&	{\hr}\tbr{t}allə	&	\cg{(olietə)}	&		&	{\hr}o\tbr{t}nə	&	{\hr}tu\tbr{t}u	&	{\hr}re\tbr{t}na	&	{\hr}mo\tbr{t}ə\\
\ili{Moa}	&	{\hr}\tbr{t}alla	&	\cg{(oilieta)}	&	{\hr}\tbr{t}eli	&	{\hr}o\tbr{t}na	&	{\hr}tu\tbr{t}u	&	{\hr}re\tbr{t}na	&	{\hr}mo\tbr{t}a\\
\ili{Leti}	&	{\hr}\tbr{t}alla	&	\cg{(olieta)}	&	{\hr}\tbr{t}eli	&	{\hr}u\tbr{t}na	&	{\hr}tu\tbr{t}u	&	{\hr}re\tbr{t}na	&	{\hr}mo\tbr{t}a\\
		\lspbottomrule
	\end{tabular}
		\begin{tablenotes}\tnsize
			\item[†]	{\ili{Roma} \it{teli} may be a loan  unless
								*z > \it{t} /{\gap}V\tsc{+fr} is a regular change.
								\ili{Kisar} \it{keleʔuk} means ‘maize, corn’.
								The origin of the final syllable \it{ʔuk} is unknown.
								\citet{Riedel1886} (in \citealp[109]{Jonker1932}) gives \ili{Kisar} \it{keli} ‘millet’.
								\ili{Moa}, \ili{Leti} \it{teli} is given as ‘maize,
								corn’ (Dutch \it{maïs}) by \citet[109]{Jonker1932}.}
			\item[‡]	{Cognates of **mozo ‘green’ occur widely in other \tsc{Timor-Babar} languages.
								Examples include: \ili{Termanu} \it{mo{\tl}modo-k} ‘yellow’,
								\ili{Meto} \it{moroʔ} ‘yellow’, \ili{Tetun} \it{modok} ‘yellow,
								green’, and Ili{\Q}uun \it{mosoŋ} ‘blue,
								green’ all of which show the expected reflexes of *z in each language.
								Putative cognates of this word also occur in languages
								of southern Sulawesi, see \srf{15-04-02-01}.}
			\item[Δ]	{Final \it{nna} in \ili{Roma} \it{atyounna} ‘far’ is analogous
								to final \it{nne} in \ili{Kisar} \it{koʔunne} ‘distance’.}
		\end{tablenotes}
	\end{threeparttable}
\end{adjustbox}
\end{table}

The strongest evidence for \tsc{\ili{Kisar}-Luang} comes from the reflexes
of *z, which are unique in the region.
\citet[289f]{Mills2010} provides a discussion of *z in these languages.
\ili{Roma} appears to have *z > \it{ty} [ʧ] before non-front vowels,
%and *z > \it{t} elsewhere (before 
and *z > \it{t} medially with raising of the following vowel to \it{i}.
The \ili{Luang} languages/dialects have *z > \it{t}
and \ili{Kisar} mergers *z/*t > \it{k}.
These changes can be accounted for by positing Proto-\ili{Kisar}-\ili{Luang}
*z > **ty [ʧ] with subsequent **ty > **t in \ili{Kisar}
and \ili{Luang}, and **t > \it{k} in \ili{Kisar}.
The reflexes of *z in are given in \trf{04-27}.
(Throughout this section \ili{Wetan} from \ili{Teun} Island [\citealp{Chlenova2006}]
is given as “\ili{Wetan} (T)”
and \ili{Wetan} from \ili{Wetan} Island
[\citealp{DeJosselinDeJong1987}] is given
as “\ili{Wetan} (W)” or just “\ili{Wetan}”.)

Another change which provides evidence for \tsc{\ili{Kisar}-Luang}
is *s > \it{h}.
This change is not complete in any language.
In particular, \ili{Leti} does not participate in this change
and \ili{Moa} only has *s > \it{h} intervocalically or before another consonant.
There are also sporadic cases of *s = \it{s} in other languages.
Examples of *s > \it{h,
s} are given in \trf{04-28}.
\ili{Wetan} has subsequent **h > {\0}
and \ili{Luang} also has **h > {\0} in at l\ili{east} one word.

\begin{table}\tcp{\tsc{\ili{Kisar}-Luang} *s}{04-28}
%\begin{adjustbox}{width=\textwidth}
	\begin{threeparttable}\st{2.5pt}
	\begin{tabular}{lllllllll}\lsptoprule
local gl.	&	br\ili{east}	&	dog	&	tooth	&	island	&	who?	&	flesh	&	hot	&	nine\\
PMP	&	*\tbr{s}u\tbr{s}u	&	*a\tbr{s}u	&	*ŋi\tbr{s}i	&	*nu\tbr{s}a	&	*\tbr{s}ai	&	*i\tbr{s}iʔ	&	*ma-pana\tbr{s}	&	*\tbr{s}iwa\\ \midrule
\ili{Roma}	&	{\hr}\tbr{h}u\tbr{h}u	&	{\hr}a\tbr{h}u	&	{\hr}ni\tbr{h}na	&	{\hr}nu\tbr{h}a	&	{\hr}\tbr{h}oi	&	{\hr}i\tbr{h}ni	&	{\hr}ma\tbr{h}an\su{†}	&	{\hr}wo-\tbr{h}ia\\
\ili{Kisar}	&	{\hr}\tbr{h}u\tbr{h}un	&	{\hr}a\tbr{h}u	&	{\hr}ni\tbr{h}in	&	{\hr}no\tbr{h}o	&	{\hr}in\tbr{h}oi	&	{\hr}i\tbr{h}i	&	{\hr}man\tbr{h}a	&	{\hr}wo-\tbr{h}ii\\
\ili{Wetan}	&	{\hr}ui	&	{\hr}ai	&	{\hr}nia	&	{\hr}noa	&	{\hr}eri(a)	&	\tcg{(pipawi)}	&	{\hr}paana\su{†}	&	{\hr}wo-\tbr{s}iwa\\
\ili{Luang}	&	{\hr}u\tbr{h}u	&	{\hr}a\tbr{h}u	&	{\hr}ni\tbr{h}ə	&	{\hr}no\tbr{h}ə	&	{\hr}\tbr{h}eʔə	&	{\hr}i\tbr{h}i	&	{\hr}pa\tbr{h}nə\su{†}	&	{\hr}wo-\tbr{s}iewə\\
\ili{Moa}	&	{\hr}\tbr{s}u\tbr{h}u	&	{\hr}a\tbr{h}u	&	{\hr}ni\tbr{h}i-ni	&	{\hr}no\tbr{h}a	&	{\hr}\tbr{s}e	&	{\hr}i\tbr{h}i	&	{\hr}pan\tbr{s}a	&	{\hr}wo-\tbr{s}iewa\\
\ili{Leti}	&	{\hr}\tbr{s}u\tbr{s}u	&	{\hr}a\tbr{s}u	&	{\hr}ni\tbr{s}a	&	{\hr}nu\tbr{s}a	&	{\hr}\tbr{s}ei	&	{\hr}i\tbr{s}i	&	{\hr}pan\tbr{s}a	&	{\hr}wo-\tbr{s}ia\\
		\lspbottomrule
	\end{tabular}
		\begin{tablenotes}\tnsize
			\item[†] {\ili{Roma}, \ili{Luang}, and \ili{Wetan} have metathesis of the final consonants from **ma{\nbhyp}pahan.}
		\end{tablenotes}
	\end{threeparttable}
%\end{adjustbox}
\end{table}

\tsc{\ili{Kisar}-Luang} languages also share *aw > \it{a} in most words.
Examples are given in \trf{04-29}.
Note that \citet[13]{TaberTaber2015} analyse \ili{Luang} word-final
[ə] as an allophone of /a/.
That is, /a/ → [ə] /{\gap\#}.

\begin{table}
	\centering\tcp{\tsc{\ili{Kisar}-Luang} *-aw}{04-29}
	\st{6pt}\begin{tabular}{llllll}\lsptoprule
local gl.	&	mouse						&	walk								&	day, sun					&	steal	&	still, calm\\
PMP				&	*balab\tbr{aw}	&	*lak\tbr{aw}				&	*qaləj\tbr{aw}		&	*nak\tbr{aw}	&	*lin\tbr{aw}\\ \midrule
\ili{Roma}			&	\cg{(melan)}		&	{\hr}na-wl\tbr{a}		&	{\hqa}ler\tbr{a}	&	{\hr}mna\tbr{a}	&	\\
\ili{Kisar}			&	\cg{(ornoho)}		&	{\hr}na-laʔ\tbr{a}	&	{\hqa}ler\cg{(e)}	&	{\hr}-amnaʔ\tbr{a}	&	{\hr}lin\tbr{a}\\
\ili{Wetan} (T)	&	\cg{(upemelai)}	&	{\hr}ta-la\tbr{a}		&	{\hqa}ler\tbr{a}	&		&	\\
\ili{Wetan} (W)	&	\cg{(upmela)}		&	{\hr}la\tbr{a}			&	{\hqa}ler\tbr{a}	&	{\hr}mna\tbr{a}	&	\\
\ili{Luang}			&	{\hr}mlaw\tbr{a}&	{\hr}na-laʔ\tbr{ə}	&	{\hqa}ler\tbr{ə}	&	{\hr}kamnaʔ\tbr{ə}	&	{\hr}lin\tbr{ə}\\
\ili{Moa}				&									&											&										&	{\hr}kamna\tbr{a}	&	\\
\ili{Leti}			&	\cg{(samɛla)}		&	{\hr}l{\tl}l\tbr{a}	&	{\hqa}lɛr\tbr{a}	&	{\hr}kamna\tbr{a}	&	{\hr}na-mnin\tbr{a}\\
		\lspbottomrule
	\end{tabular}
\end{table}

\newpage
\tsc{\ili{Kisar}-Luang} languages also share *j > \it{r} in most words,
with examples given in \trf{04-30}.
However, reflexes of *huaji ‘younger sibling’ in \ili{Roma}
and \ili{Kisar} have *j > \it{ll},
reflexes of *ŋajan ‘name’ have *j > {\0} in the \ili{Luang} language/dialect
cluster, and the reflex of *j in *laləj is ambiguous between \it{r} or
\it{n} in the \ili{Luang} language/dialect cluster.

\begin{table}\tcp{\tsc{\ili{Kisar}-Luang} *j}{04-30}
%\begin{adjustbox}{width=\textwidth}
	\begin{threeparttable}\st{2.75pt}
	\begin{tabular}{llllllll}\lsptoprule
local gl.	&	sun	&	how m.?	&	nose	&	ySib\su{†}	&	name	&	fly	&	navel\\
PMP	&	*qalə\tbr{j}aw	&	*pi\tbr{j}a	&	*i\tbr{j}uŋ	&	*hua\tbr{j}i	&	*ŋa\tbr{j}an	&	*lalə\tbr{j}	&	*pusə\tbr{j}\\	\midrule
\ili{Roma}	&	{\hqa}le\tbr{r}a	&	{\hr}wo-i\tbr{r}a	&	{\hr}i\tbr{r}un	&	{\hr}a\tbr{ll}i	&	{\hr}na\tbr{r}an-na	&	{\hr}lara\tbr{r},	&	{\hr}uha\tbr{r}\\
			&	&	&	&	&	&	{\hr}lar\tbr{r}a	&	\\
\ili{Kisar}	&	{\hqa}le\tbr{r}	&	{\hr}wo-i\tbr{r}a	&	{\hr}i\tbr{r}un	&	{\hr}a\tbr{ll}a	&	{\hr}na\tbr{r}an	&	{\hr}lar\tbr{r}a	&	{\hr}oho\tbr{r}-ne\\
\ili{Wetan}	&	{\hqa}le\tbr{r}a	&	{\hr}wo-i\tbr{r}a	&	{\hr}i\tbr{r}na	&	{\hr}ya\tbr{r}kiea	&	{\hr}nana	&		&	{\hr}o\tbr{r}a\\
\ili{Luang}	&	{\hqa}le\tbr{r}ə	&	{\hr}wo-i\tbr{r}ə	&	{\hr}i\tbr{r}nu	&	{\hr}ga\tbr{r}i	&	{\hr}nanə	&	{\hr}larna	&	{\hr}goho\tbr{r}ə\\
\ili{Moa}	&	{\hqa}le\tbr{r}a	&	{\hr}wo-i\tbr{r}a	&	{\hr}i\tbr{r}nu-e	&	{\hr}ga\tbr{r}i	&	{\hr}nana	&	{\hr}llarna	&	{\hr}oha\tbr{r}-ni\\
\ili{Leti}	&	{\hqa}le\tbr{r}a	&	{\hr}βo-i\tbr{r}a	&	{\hr}i\tbr{r}nu	&	{\hr}wa\tbr{r}i	&	{\hr}naana	&	{\hr}llarna	&	{\hr}oh\tbr{r}a\\
		\lspbottomrule
	\end{tabular}
		\begin{tablenotes}\tnsize
			\item[†] {\ili{Kisar} \it{alla} is only currently known
								in the phrase \it{kaka-alla} ‘relatives’
								(\it{kaka} = ‘older relatives’).
								\ili{Wetan} \it{yarkiea} is from \it{yari + kea}.
								\it{yari} is given as ‘last, youngest´.
								Other reflexes are given as ‘younger sibling’.}
		\end{tablenotes}
	\end{threeparttable}
%\end{adjustbox}
\end{table}

\tsc{\ili{Kisar}-Luang} languages also have *b > **β in most cases.
The fricative realisation [β] is retained in \ili{Leti},
while other descriptions tend to describe it as a glide.
\ili{Roma} and \ili{Luang} /w/ is described as “[{\ldots}] a voiced bilabial glide.
[{\ldots}] with less lip rounding
and more friction than the \ili{English} [w]”
(\citealp[9]{TaberTaber2015}, \citealp[34]{Steven1991}),
with allophones [ʋ] and [v] depending on adjacent vowels.
Reflexes are given in \trf{04-31}.
\ili{Roma} and \ili{Kisar} have subsequent **β > {\0} /V{\gap}u.

\begin{table}\tcp{\tsc{\ili{Kisar}-Luang} *b}{04-31}
%\begin{adjustbox}{width=\textwidth}
	\st{5.5pt}\begin{tabular}{lllllll}\lsptoprule
local gl.	&	fruit	&	moon	&	stone	&	pig	&	ash, dust	&	sugarcane\\
PMP	&	*\tbr{b}uaq	&	*\tbr{b}ulan	&	*\tbr{b}atu	&	*\tbr{b}a\tbr{b}uy	&	*qa\tbr{b}u	&	*tə\tbr{b}uh\\ \midrule
\ili{Roma}	&	{\hr}\tbr{w}una	&	{\hr}\tbr{w}ulan	&	{\hr}\tbr{w}atu	&	{\hr}\tbr{w}a\tbr{w}i	&	{\hr}au	&	{\hr}teu\\
\ili{Kisar}	&	{\hr}\tbr{w}oi	&	{\hr}\tbr{w}ollo	&	{\hr}\tbr{w}aku	&	{\hr}\tbr{w}a\tbr{w}i	&	{\hr}a\tbr{p}u	&	{\hr}keu\\
\ili{Wetan} (T)	&	{\hr}\tbr{w}oa	&	{\hr}\tbr{w}olna	&	{\hr}\tbr{w}ati	&	{\hr}\tbr{w}a\tbr{w}i	&	{\hr}a\tbr{w}i	&	\\
\ili{Wetan} (W)	&	{\hr}\tbr{w}oa	&	{\hr}\tbr{w}olla	&	{\hr}\tbr{w}ati	&	{\hr}\tbr{w}a\tbr{w}i	&	{\hr}a\tbr{w}i	&	{\hr}te\tbr{w}i\\
\ili{Luang}	&	{\hr}\tbr{w}oʔə	&	{\hr}\tbr{w}ollə	&	{\hr}\tbr{w}atu	&	{\hr}\tbr{w}a\tbr{w}i	&	{\hr}a\tbr{w}u	&	{\hr}te\tbr{w}u\\
\ili{Moa}	&	{\hr}\tbr{w}oa	&	{\hr}\tbr{w}olla	&	{\hr}\tbr{w}atg-e	&	{\hr}\tbr{w}a\tbr{w}i	&	{\hr}a\tbr{w}g-e	&	\\
\ili{Leti}	&	{\hr}\tbr{β}ua, \tbr{β}ɔa	&	{\hr}\tbr{β}ulna	&	{\hr}\tbr{β}atu	&	{\hr}\tbr{β}a\tbr{β}i	&	{\hr}a\tbr{β}u	&	{\hr}te\tbr{β}u\\
		\lspbottomrule
	\end{tabular}
%\end{adjustbox}
\end{table}

\tsc{\ili{Kisar}-Luang} also has *ŋ/*n > \it{n} in most positions.
\citet[28ff]{Mills2010} shows that complete merger of *ŋ/*n > \it{n}
in all word positions cannot be posited at the level of Proto-\ili{Kisar}-\ili{Luang}
as the \ili{Luang} languages/dialects show *ŋ > \it{k} /{\gap}\it{r}.
Examples are shown in \trf{04-32}.

\begin{table}\tcp{\tsc{\ili{Kisar}-Luang} *ŋ}{04-32}
%\begin{adjustbox}{width=\textwidth}
	\begin{threeparttable}\st{4pt}
	\begin{tabular}{lllllll}\lsptoprule
local gl.	&	name	&	tooth	&	wind	&	ear	&	soul	&	fry\\
PMP	&	*\tbr{ŋ}ajan	&	*\tbr{ŋ}isi	&	*ha\tbr{ŋ}in	&	*tali\tbr{ŋ}a\su{†}	&	*suma\tbr{ŋ}əd	&	*sa\tbr{ŋ}əlaR\su{†}\\ \midrule
\ili{Roma}	&	{\hr}\tbr{n}aran-na	&	{\hr}\tbr{n}ihna	&	{\hr}a\tbr{n}ni	&	{\hr}tli\tbr{n}-na	&	{\hr}hamara\tbr{n}-na\su{‡}	&	{\hr}hora\tbr{n}\su{‡}\\
\ili{Kisar}	&	{\hr}\tbr{n}aran	&	{\hr}\tbr{n}ihin	&	{\hr}a\tbr{n}ne	&	{\hr}kili\tbr{n}-na	&	{\hr}hamara\tbr{n}\su{‡}	&	{\hr}har\tbr{n}a\su{‡}\\
\ili{Wetan}	&	{\hr}\tbr{n}ana	&	{\hr}\tbr{n}ia	&	{\hr}a\tbr{n}ni	&	{\hr}tli\tbr{n}a	&	{\hr}ma\tbr{k}ra	&	\\
\ili{Luang}	&	{\hr}\tbr{n}anə	&	{\hr}\tbr{n}ihə	&	{\hr}a\tbr{n}ni	&	{\hr}tli\tbr{n}ə	&	{\hr}(h)ma\tbr{k}ar\su{Δ}	&	{\hr}he\tbr{k}rə\\
\ili{Moa}	&	{\hr}\tbr{n}ana	&	{\hr}\tbr{n}ihi-ni	&	{\hr}a\tbr{n}ni	&	{\hr}tli\tbr{n}a	&		&	\\
\ili{Leti}	&	{\hr}\tbr{n}aana	&	{\hr}\tbr{n}isa	&	{\hr}a\tbr{n}ni	&	{\hr}tni\tbr{n}a	&	{\hr}sma\tbr{k}ra	&	{\hr}sɛ\tbr{k}ra\\
		\lspbottomrule
	\end{tabular}
		\begin{tablenotes}\tnsize
			\item[†]	{*saŋəlaR is reconstructed with the gloss ‘stir-fry,
								cook in a frying pan without oil’.
								The words in this table are inherited from intermediate **səŋaR
								(also attested in other \tsc{Timor-Babar} languages,
								see \trf{04-12} on \prf{tab:04-12}).
								\citet{BlustTrussel2020} provide a helpful summary of the phonological
								changes that have affected this word.
								\ili{Kisar} \it{harna} means ‘burn, roast’.}
			\item[‡]	{\ili{Roma} and \ili{Kisar} reflexes of *sumaŋəd
								and *saŋəlaR have metathesis of **nVr > \it{rVn}.}
			\item[Δ]	{\ili{Luang} \it{(h)makar} is from \citet[192]{Jonker1932}
								and represents the Sermata variety of \ili{Luang}.}
		\end{tablenotes}
	\end{threeparttable}
%\end{adjustbox}
\end{table}

Within the \tsc{\ili{Kisar}-Luang} node,
\ili{Kisar} and \ili{Roma} can be placed together based on shared *wa > \it{o}
in reflexes of *wahiyəR ‘water’
and *w > {\0} in other identified reflexes,
though \ili{Kisar} retains *w in *wani > \it{wani} ‘bee’.
Additional evidence for \tsc{\ili{Kisar}-Roma} comes from shared *ma-p
> \it{m} in reflexes of *ma-panas ‘hot’ as opposed to *ma-p >
\it{p} in the other languages (see \trf{04-28})
and *j > \it{ll} (rather than *j  > \it{r}) in reflexes of *huaji ‘younger sibling’.
Against this evidence for \tsc{\ili{Kisar}-Roma},
is shared *ay > \it{i} in \ili{Kisar}
and the \ili{Luang} languages/dialects,
while \ili{Roma} has *ay > \it{a}.
Examples of the reflexes of *w
and *ay in the \tsc{\ili{Kisar}-Luang} languages are shown in \trf{04-33}.

\clearpage
\begin{table}\tcp{\tsc{\ili{Kisar}-Luang} *w and *-ay}{04-33}
\begin{adjustbox}{width=\textwidth}
	\st{2pt}\begin{tabular}{lll@{}ll|llll}\lsptoprule
local gl.	&	water	&	root	&	ySib	&	eight	&	liver	&	die	&	sand	&	paddle\\
PMP	&	*\tbr{w}ahiyəR	&	*\tbr{w}akaR	&	*h\tbr{ua}ji	&	*\tbr{w}alu	&	*qat\tbr{ay}	&	*mat\tbr{ay}	&	*qən\tbr{ay}	&	*bəRs\tbr{ay}\\ \midrule
\ili{Roma}	&	{\hr}\tbr{o}ri	&	\cg{(muar-na)}	&	{\hr}alli	&	{\hr}wo-allu	&	{\hr}at\tbr{a}n-na	&	\cg{(m-pora)}	&	{\hr}en\tbr{a}	&	{\hr}woh\tbr{a}\\
\ili{Kisar}	&	{\hr}\tbr{o}ir	&	{\hr}aʔarn	&	{\hr}alla	&	\cg{(wo-aa)}	&	{\hr}ak\tbr{i}n	&	{\hr}mak\tbr{i}	&	\cg{(ʈosle)}	&	{\hr}woh\tbr{i}\\
\ili{Wetan}	&	{\hr}\tbr{y}era	&	{\hr}\tbr{y}aara	&	{\hr}\tbr{y}ari	&	\cg{(wo-awa)}	&	{\hr}at\tbr{i}	&	{\hr}mat\tbr{i}	&	{\hr}en\tbr{i}	&	{\hr}we\tbr{i}\\
\ili{Luang}	&	{\hr}\tbr{g}erə	&	{\hr}\tbr{g}aʔrə	&	{\hr}\tbr{g}ari	&	\cg{(wo-awə)}	&	{\hr}at\tbr{i}	&	{\hr}mat\tbr{i}	&	{\hr}en\tbr{i}	&	\\
\ili{Moa}	&	{\hr}\tbr{g}era	&	{\hr}\tbr{g}aara	&	{\hr}\tbr{g}ari	&	\cg{(wo-awa)}	&	{\hr}at\tbr{i}	&	{\hr}mat\tbr{i}	&	{\hr}en\tbr{i}	&	{\hr}weh\tbr{i}\\
\ili{Leti}	&	{\hr}\tbr{w}ɛra	&	{\hr}\tbr{w}aara	&	{\hr}\tbr{w}ari	&	\cg{(βɔ-aβa)}	&	{\hr}at\tbr{i}	&	{\hr}mat\tbr{i}	&	{\hr}en\tbr{i}	&	{\hr}βes\tbr{i}\\
		\lspbottomrule
	\end{tabular}
\end{adjustbox}
\end{table}
